\section{Composing: Paulis, the Semidirect Product, and Scalars}\label{sec:composing}

$\PPauli \cong (\Zp \times \Zp)^n$ is presented over $X$, $Z$ on each
wire by $X^{p} = Z^{p} = \varepsilon$ and $ZX = XZ$; its semantic
group is the phase space that $\Sp$ acts on, and its normal form is
$X^{a}Z^{b}$ per wire.  Gluing this to the symplectic presentation of
\S\ref{sec:symplectic} is the semidirect example of
Appendix~\ref{sec:sn-wreath} with $\Sp$ for $S_n$ and $\PPauli$ for
$(\mathbb{Z}_N)^n$: the presentation has both factors' relations plus
the conjugation relations $h\,n\,h^{-1} = \mathit{conj}\,h\,n$, the
action \aid{conj} spelling out the paper's Lemmas~2.24--2.25.  The
theorem's two conditions --- the action respects the relations of
each factor --- are here proved by semantics: the conjugates of
congruent words have equal denotations by symplectic soundness, hence
are congruent by Pauli \emph{completeness} (\aid{respects-Δ}; the
paper's ``Pauli complete'' rules, Definition~4.7).  Every premise
being a theorem, the semidirect presentation of \aid{Pauli⋊Sp} is
unconditional.

\emph{Plumbing.}  The semidirect presentation has generators $X, Z, H,
S, CZ$ and $3 + 15 + 14$ relations (Pauli, symplectic, conjugation);
the paper's rule sets have $H, S, CZ$ alone.  Three isomorphisms of
presentations close the gap (Figure~\ref{fig:route}), through
\texttt{Simplified-V1} (nineteen rules), \texttt{Paper-V0} (sixteen)
and \texttt{Paper-V1} (Theorem~\ref{thm:main}'s
fifteen): generator reduction ($X$ and $Z$ become words) and relation
reduction (the Pauli and conjugation rules become theorems).
\emph{Semidirect to Simplified-V1} is an isomorphism on \emph{different}
alphabets, $f$ sending $S$ to $R = S \bullet Z^{1/2}$ and $h$ sending
$S$ to $Z^{-1/2} \bullet S$; in particular, $f$ induces the embedding
of the symplectic group into the projective Clifford
group.  \texttt{Simplified-V1}'s nineteen rules are the fifteen simplified symplectic ones
plus $(SH)^3 = \varepsilon$, $H^2SH^2S = SH^2SH^2$ and the two rules
pushing $X$ through $CZ$; \aid{f-well-defined} derives from them the
images under $f$ of the three Pauli and the fourteen conjugation
rules (6\,000 lines).  \emph{Simplified-V1 to Paper-V0} is
the identity on circuits: ten axioms are shared, and V0's \texttt{Lemmas}
(4\,750 lines) derives c10--c15 from its own axioms.
\emph{Paper-V0 to Paper-V1} reaches the paper's own rules.  The two
rule sets differ in how the multiplier is spelled as a word:
\texttt{Paper-V0} spells $M_x = R^{x}HR^{x^{-1}}HR^{x}H$ over $R$;
the paper's~(T1) spells $M_x = \aid{SHS'}\;x^{-1}\;x$ over $S$, with
$\aid{SHS'}\;a\;b = Z^{(b-1)/2}X^{(1-a)/2}S^{b}HS^{a}HS^{b}H$ and
\aid{semi-MR} as $M_gS^{g^2}Z^{(g^2-g)/2} = SM_g$.  \texttt{Paper-V1}
restates the four multiplier rules over \aid{SHS'} and drops one:
$M_1$ is now the word $(SH)^3$, so \aid{M-power} at $k = 0$ already
gives $(SH)^3 \approx \varepsilon$, whereas over $R$ it is $(RH)^3$,
a different word, and the axiom had to stay.  The bridge
$R^{a}HR^{b}HR^{a}H \approx \aid{SHS'}\;a\;b$
(\texttt{Shared.PauliBase}) carries the multiplier rules across,
giving Theorem~\ref{thm:main}.

\emph{Scalars.}  \aid{extension-presentation} extends the projective
presentation --- \texttt{Paper-V0}'s, which has \aid{order-SH} as an
axiom --- by the scalar (Figure~\ref{fig:exact}): add $\omega$ and $\omega^{p} = \varepsilon$,
let $\omega$ commute with every gate, and \emph{correct} each rule by
its scalar.  The correction is $\varepsilon$ except on
\aid{order-SH}, the one rule failing on the nose: $(SH)^3 =
\omega^{-1/8}$, exponent $(p^2{-}1)/8 \equiv -1/8$ in $\Zp$.
The multiplier words now denote $\ket{j} \mapsto \ket{xj}$
only up to a scalar.  The semantic extension is
the Weyl cocycle on the natural semidirect product, $c((P,S),(Q,S'))
= \tfrac12\,\mathrm{sform}(P,\ S \cdot Q)$, with $S \cdot Q$ the
symplectic part acting on the Pauli vector and $\mathrm{sform}$ the
symplectic form of the phase space, the exponent in $PQ =
\omega^{\mathrm{sform}(P,Q)}\,QP$; its cocycle identity is, by
bilinearity, symplectic invariance of \aid{sform}, and the
presented group's cocycle is the pullback of $c$ along
Theorem~\ref{thm:main}'s isomorphism, so the two agree by
construction.  The unitary identification uses Weyl multiplication
\cite[Eq.~(4), Theorem~3]{gross2006hudson} and a \emph{linear} Weil
representation with exact covariance
\cite[Theorem~1.3.2]{gurevich2007geometric}: the operators
$U(e,v,S)=\omega^e W(v)\rho(S)$ obey
\[
 U(e,v,S)U(f,w,T)
 =U\bigl(e+f+\tfrac12\mathrm{sform}(v,Sw),v+Sw,ST\bigr).
\]
This is the Heisenberg--Weil representation
\cite[\S 1.3]{gurevich2010bounds}; conjugation on Weyl operators
determines $S$ and $v$, so it is faithful and its scalars are exactly
$\langle\omega I\rangle$ (at width zero, $e\mapsto\omega^e$).
Adjoining $-I$ gives its direct product with $\mathbb Z_2$ and scalars
$\langle-\omega I\rangle$.  These finite phase choices differ from
the full unitary normalizer, which contains $U(1)$.
Unconditionally:

\begin{code}%
\>[0]\AgdaFunction{presentation-exact}\AgdaSpace{}%
\AgdaSymbol{:}\AgdaSpace{}%
\AgdaSymbol{∀}\AgdaSpace{}%
\AgdaSymbol{(}\AgdaBound{n}\AgdaSpace{}%
\AgdaSymbol{:}\AgdaSpace{}%
\AgdaDatatype{ℕ}\AgdaSymbol{)}\AgdaSpace{}%
\AgdaSymbol{→}\<%
\\
\>[0][@{}l@{\AgdaIndent{0}}]%
\>[2]\AgdaKeyword{let}\AgdaSpace{}%
\AgdaBound{Clifford-group}\AgdaSpace{}%
\AgdaSymbol{=}\AgdaSpace{}%
\AgdaSymbol{λ}\AgdaSpace{}%
\AgdaBound{n}\AgdaSpace{}%
\AgdaSymbol{→}\AgdaSpace{}%
\AgdaPostulate{PE.Presented-group}\AgdaSpace{}%
\AgdaBound{n}\AgdaSpace{}%
\AgdaSymbol{(}\AgdaPostulate{exact-generator-data}\AgdaSpace{}%
\AgdaBound{n}\AgdaSymbol{)}\<%
\\
%
\>[2]\AgdaKeyword{in}\AgdaSpace{}%
\AgdaSymbol{(}\AgdaBound{n}\AgdaSpace{}%
\AgdaOperator{\AgdaPostulate{Exact,\AgdaUnderscore{}===\AgdaUnderscore{}}}\AgdaSymbol{)}\AgdaSpace{}%
\AgdaOperator{\AgdaRecord{IsPresentationOf}}\AgdaSpace{}%
\AgdaSymbol{(}\AgdaBound{Clifford-group}\AgdaSpace{}%
\AgdaBound{n}\AgdaSymbol{)}\<%
\end{code}


\noindent The extension is provably non-split.

\emph{Figure~1 as printed} is over $-\omega$, its multipliers and swap
carrying Legendre symbols and $\lambda_p^2$, so it needs the $-1$ that
the exact group (scalars of odd order) lacks.  Adjoining $-1$ as a
central involution presents the exact group times $\mathbb Z_2$,
\aid{Clifford±}, and Figure~1 is isomorphic to that rule set by
$-\omega \mapsto -1\cdot\omega$, each exact rule's scalar pinned by
reading its lift back: the paper's Theorem~4.10, against our group,
%% Copied from ../vqupit/latex/vqupit/sections/composing.tex (composing, the figure1-presentation block), the LaTeX
%% that `agda --latex --only-scope-checking' generated for the long version.
\begin{code}%
\>[0]\AgdaFunction{figure1-presentation}\AgdaSpace{}%
\AgdaSymbol{:}\AgdaSpace{}%
\AgdaSymbol{∀}\AgdaSpace{}%
\AgdaBound{n}\AgdaSpace{}%
\AgdaSymbol{→}\AgdaSpace{}%
\AgdaSymbol{(}\AgdaBound{n}\AgdaSpace{}%
\AgdaOperator{\AgdaPostulate{F,\AgdaUnderscore{}===\AgdaUnderscore{}}}\AgdaSymbol{)}\AgdaSpace{}%
\AgdaOperator{\AgdaRecord{IsPresentationOf}}\AgdaSpace{}%
\AgdaSymbol{(}\AgdaPostulate{Clifford±}\AgdaSpace{}%
\AgdaBound{n}\AgdaSymbol{)}\<%
\end{code}

